\section{Conclusions}
\label{sec:conclusions}

This paper asked whether any of the nearby stars that ALMA has observed shows
a narrowband, drifting radio carrier of the kind an artificial transmitter
produces at millimetre and submillimetre wavelengths, and, if none is
seen, how deep and how complete a search of the archive for such a carrier
can be. To answer it we searched the archival ALMA spectra of
\NSysClassA{} star systems within 40\,pc. The fine-channel windows cover
\CvUnionA\,GHz of sky frequency, in \CovNIslands{} separate intervals between
\SurvFreqLoA{} and \SurvFreqHiA\,GHz, which is at least \SbPriorRatio{} times
the bandwidth of either of the two millimetre searches of nearby stars that
preceded it. We removed the drift from
each spectrum and compared it with control positions in the same field, so
that the background at the star could be judged from the observation itself.

No transmitter was found. In \CsNCross{} windows the stellar statistic reaches
the threshold $T_\star\geq\StageNullTrig$. Of these, \MkNCoinc{} coincide with
catalogued molecular lines, \LnAttrSupp{} of them with further evidence of
line emission, and none was identified as interference by the available
diagnostics. The other \CeObs{} are consistent with chance: the
control positions beside each star predict \CeExp, with $P=\CeProbNull$
(\FrTwoTrialP{} for the alternative extraction field after allowing for the
choice, neither significant at the \CeAlphaPct{} per cent level). That
probability is conditional on the line mask as adopted, which was completed
after the crossings had been examined, and on the control positions being a
fair model of the background at the star. Wherever an unattributed crossing
was observed again and the repeat could be read, it did not return at its
predicted frequency, so no recurrence has been confirmed. This is a
non-detection, not an exclusion.

Because nothing was detected, what the search delivers is a set of upper
limits. They apply to a carrier narrower than one channel, lying outside the
molecular-line mask, radiating while ALMA observed and drifting by no more
than \DriftCeilLo--\DriftCeilHi\,nHz, the drift ceiling discussed in
\S\ref{sec:excludes}. For such a carrier, recovery nine times in ten in a single epoch needs an equivalent
isotropic radiated power, or EIRP, with a median, over systems, of
$\EirpNinetySysMedA$\,W in each system's most sensitive window. The deepest
limit, $\EirpNinetySysLoA$\,W, is reached toward $\epsilon$~Eridani, and
\NSysEfifteen{} of the \NSysALimSum{} systems whose limits do not depend on an
unchecked position reach $\CovEirpRefTex$\,W.

It helps to set these numbers against something we have built ourselves.
The Arecibo planetary radar, one of the most powerful transmitters ever
operated on the Earth, had an EIRP of about $\AroEirp$\,W at
\AroFreqGHz\,GHz. The median limit found here is about \AroMedMult{} times
that, and even the deepest is \AroDeepMult{} times it, so
\AroNSysBelowWord{} system reaches an Arecibo-like transmitter, although
\AroNSysTen{} come within a factor of ten of it. The same comparison can be
put as a distance. Because the received flux falls as the square of the
distance, an Arecibo-like transmitter placed at a star of this survey would
be detected only if the system were much closer to us than it is. For the
median system it would have to be within \AroDistMedPc\,pc, and even in the
most favourable case, \AroDistMaxStar, whose limit is deep for its distance,
within \AroDistMaxPc\,pc. A \BenchDiam\,m dish radiating \BenchPowerMW\,MW is
a more natural transmitter at these frequencies. Tuned to the frequency of
each system's most sensitive window, between \BenchTunedFreqLo{} and
\BenchTunedFreqHi\,GHz, its EIRP lies above the single-epoch limit of
\BenchNSysTuned{} of the \NSysALimSum{} systems. It could only be seen, of
course, if its beam happened to point at us, as discussed in
\S\ref{sec:mmdifferent}.

How complete the search can be is set more by the archive's coverage than by
its depth. A crossing becomes a testable claim only when its frequency is
observed again in a separated session. Of the \NSysClassA{} systems,
\ZcNoRepNoSep{} were never observed in two separated sessions at all, so that
nothing found toward them could have been confirmed at any power. For
\ZcNoRepNoFreqWord{} more, separated sessions exist but share no covered
frequency, so the same-frequency recurrence test cannot be applied there. A
programme designed for the purpose, on the same telescope, would use fine
channels throughout, tune away from the bright molecular lines and schedule its repeat observations deliberately, so that
any crossing it found could be looked for again at the frequency it predicts
(\S\ref{sec:future}).
